% Public-safe skill scan appendix.
% This is a narrative appendix for the public research-status report. It does
% not render internal STEM scores, trust scores, repository lists, restricted
% project names, or private evidence.

\section*{Reported Skill Scan Appendix}

This section summarizes a public-safe skill scan of repository-derived work. The scan is used as a documentation aid: it connects repeated work patterns to visible skill demonstrations, separates mature evidence from developing areas, and avoids presenting internal scores as credentials.

\subsection*{Market-Ready Skill Demonstrations}

\begin{itemize}
    \item \textbf{R analytics:} Demonstrated through longitudinal research workflows, statistical modeling, table generation, sample-audit discipline, and manuscript-oriented analysis. The evidence supports applied research roles where analytic outputs must be traceable from source data to written claim.
    \item \textbf{Python:} Demonstrated through build scripts, validation checks, object generation, public-release guards, manifest contracts, and automation for repeatable documentation. The evidence supports scripting, data preparation, workflow validation, and reproducible reporting.
    \item \textbf{Data engineering:} Demonstrated through structured manifests, source ledgers, generated project objects, release-control records, and public/private boundary handling. The evidence supports analytic dataset preparation, evidence organization, and controlled publication pipelines.
    \item \textbf{Cloud and DevOps workflows:} Demonstrated through GitHub Actions, Makefile-driven builds, artifact packaging, contract checks, and reproducible public-output generation. The evidence supports CI-based documentation systems and reviewable build automation.
    \item \textbf{Cybersecurity-aware data governance:} Demonstrated through public-release filtering, index-safe output controls, hashed release guards, claim-boundary checks, and separation of public artifacts from private source records. The evidence supports careful data handling, auditability, and public-safe publishing.
    \item \textbf{AI-assisted evidence engineering:} Demonstrated through role-object packaging, scanner interpretation, review gates, and evidence-alignment workflows. The evidence supports governed use of automation for documentation, application support, and quality-control review rather than unsupported claims of model-platform ownership.
\end{itemize}

\subsection*{Skills in Training}

\begin{itemize}
    \item \textbf{Machine-learning evaluation:} Active development area supported by benchmark planning, model-card thinking, responsible evaluation notes, and early biomedical/noise-analysis scaffolds. This is treated as training-stage unless a role calls specifically for evaluation literacy rather than mature production ownership.
    \item \textbf{Geospatial and environmental-health analysis:} Active development area supported by public-health infrastructure, WASH, terrain/access, and environmental screening work. This is treated as a supporting specialization while validation records and official-source confirmation continue to mature.
\end{itemize}

\subsection*{Use Boundary}

The skill scan is not an external credential, salary estimate, or ranking system. It is a public-safe summary of demonstrated work patterns. Mature skills are listed only where repeated project evidence supports them; developing skills are labeled as training-stage rather than presented as market-ready capabilities.
